\begin{helper}[Inequality (3.3)]
Let $n\ge2$, $p>1$ and $r>0$.  Then there is a constant
$c=c\bigl(n,p,r\bigr)>0$ with the following property.  Let
$\varepsilon>0$, $\gamma_0>0$, $\gamma_1\le r\gamma_0$ and $t>0$, let
$\Omega\subseteq\R^n$ be open, let $A$ satisfy (Q1) of the thesis--(Q2) of the thesis on
$\Omega$ with parameters $p,\varepsilon,\gamma_0,\gamma_1$
(\noderef{EllipticStructureQ12}), assume the cited regularity result (C1) for
$\varepsilon>0$ (\noderef{CitedRegularityEpsPos}) with these parameters on
$\Omega$ for the field $A$, let $u$ be a weak solution of
$\operatorname{div}A(\nabla u)=0$ in $\Omega$ (\noderef{IsWeakSolutionOn}) with
weak gradient $g=\nabla u$ and with $u\in W^{2,2}_{\mathrm{loc}}(\Omega)$, the
weak gradient of $g^k=D_ku$ being $Dg^k$ (\noderef{MemW1pLocOn}).  Put
$w=\varepsilon+|\nabla u|^2$.  Then for every $\zeta\in C_0^\infty(\Omega)$,
\[
\int_\Omega\Bigl|\nabla\Bigl(w^{\frac{p+t}4}\Bigr)\Bigr|^2\zeta^2\,dx
\;\le\;c\int_\Omega w^{\frac{p+t}2}|\nabla\zeta|^2\,dx .
\]
Here the left-hand side is written out using
$\nabla(w^{s})=s\,w^{s-1}\nabla w$ and
$D_iw=2\sum_k D_ku\,D_iD_ku$, i.e.\ with $s=(p+t)/4$,
\[
\Bigl|\nabla\bigl(w^{s}\bigr)\Bigr|^{2}
=s^{2}\,w^{2s-2}\sum_{i=1}^n\Bigl(2\sum_{k=1}^nD_ku\,D_iD_ku\Bigr)^{2},
\qquad 2s-2=\frac{p+t}2-2 ,
\]
and $|\nabla\zeta|^2=\sum_i(D_i\zeta)^2$.

The hypothesis (C1) is not an addition to the thesis' setting but a transcription
of it.  The thesis derives (3.3) of the thesis by substituting a test function into
the differentiated equation (3.2) of the thesis, and it is at
(3.2) of the thesis that the regularity of $u$ enters: the thesis writes ``From now on
we assume $\varepsilon>0$, so that weak solutions of equation (3.1) of the thesis are
in the class $C^2$'' (p.~14) and the Note for the formalization replaces that
$C^2$ by exactly (C1), i.e.\ by $u\in W^{2,2}_{\mathrm{loc}}\cap C^1$.
\noderef{DifferentiatedEquationWeak} therefore carries (C1) as an explicit
hypothesis, and so must any node whose proof substitutes into it;
$\varepsilon>0$ is in force here for the same reason.  Both consumers of this
node assume (C1) in the universally quantified form the paper's Note states it
in, so the hypothesis is discharged where it is used and nothing downstream has
to be weakened.
\end{helper}
\begin{proof}
Throughout, $D_i$ is the derivative in the direction of the $i$-th standard
basis vector, $(\cdot,\cdot)$ and $|\cdot|$ are the Euclidean inner product and
norm, indices $i,j,k$ run over $\{1,\dots,n\}$, and all non-integer powers are
real powers of positive numbers.  Write
\[
w=\varepsilon+|g|^2,\qquad
M^k=Dg^k,\qquad
\nabla w=2\sum_{k=1}^ng^kM^k,\qquad
a_{ij}(x)=D_{h_j}A^i\bigl(g(x)\bigr),
\]
so that $(M^k)^j=D_jD_ku$ and $(\nabla w)_i=2\sum_kD_ku\,D_iD_ku$, which is the
field appearing in the statement; $\nabla w$ and $a_{ij}$ are notation for those
expressions and no differentiability of $w$ or of $A$ beyond what is assumed is
implied by writing them.  Finally put
\[
s=\frac{p+t}4>0,\qquad
I=\int_\Omega w^{\frac{p+t}2-2}|\nabla w|^2\zeta^2\,dx,\qquad
J=\int_\Omega w^{\frac{p+t}2}|\nabla\zeta|^2\,dx ,
\]
so that the inequality to be proved is exactly $s^2I\le cJ$, the left-hand side
of the statement being $s^2w^{2s-2}|\nabla w|^2\zeta^2$ integrated and
$2s-2=\frac{p+t}2-2$.

\emph{Step 0: the constant, and the degenerate case.}  Put
\[
c=n^2p^2r^2 ,
\]
which is positive and depends only on $n$, $p$ and $r$, and which is chosen
before $\varepsilon,\gamma_0,\gamma_1,t,\Omega,A,u,g,Dg,\zeta$ --- the
quantifier order the statement prescribes, and the formal content of
``$c=c(n,p,\gamma_1/\gamma_0)$'', the structure constants entering only through
the bound $\gamma_1\le r\gamma_0$.  If $\Omega=\emptyset$ both integrals are
over the empty set, hence $0$, and the inequality reads $0\le c\cdot0$; so
assume $\Omega\ne\emptyset$ from now on.  Being open and nonempty, $\Omega$ then
has positive measure.

\emph{Step 1: $\gamma_0\le\gamma_1$.}  This is needed because
\noderef{DifferentiatedEquationWeak} assumes it.  By
\noderef{EllipticStructureQ12} the conditions (Q1) of the thesis and (Q2) of the thesis
hold for almost every $x\in\Omega$, so --- $\Omega$ having positive measure ---
there is a single $x\in\Omega$ at which both hold.  Fix such an $x$, take $z=0$,
$h=e_1$ and $\xi=e_1$; the guard $\varepsilon+|h|^2\ne0$ of the two conditions
holds because $\varepsilon>0$.  Then
\[
\gamma_0(\varepsilon+1)^{\frac p2-1}
\le\bigl(D_hA(h)e_1,e_1\bigr)=\bigl(D_hA(h)e_1\bigr)^1
\le\bigl|\bigl(D_hA(h)e_1\bigr)^1\bigr|
\le\gamma_1(\varepsilon+1)^{\frac p2-1},
\]
and $(\varepsilon+1)^{p/2-1}>0$, so $\gamma_0\le\gamma_1$.

\emph{Step 2: what (C1) supplies, and the elementary properties of $w$.}
Instantiate \noderef{CitedRegularityEpsPos} with the present
$n,p,\varepsilon,\gamma_0,\gamma_1,\Omega,A$: its antecedents are
$\varepsilon>0$, the structure conditions on $\Omega$
(\noderef{EllipticStructureQ12}) and the weak-solution property
(\noderef{IsWeakSolutionOn}), all of which are hypotheses here, and among its
conclusions is a field $g'$ with
\[
g'\ \text{continuous on }\Omega,
\qquad
g=g'\ \text{ almost everywhere on }\Omega .
\tag{1}
\]
\emph{It does not assert continuity of the given $g$, and cannot: a weak
gradient is an almost-everywhere class, so continuity is a property of a
representative.  Every \emph{pointwise} consequence below is therefore drawn for
the primed quantities built from $g'$, and transported back to the unprimed ones
across $(1)$, which every integral and every almost-everywhere assertion is
blind to.}  Accordingly put
\[
w'=\varepsilon+|g'|^2,\qquad
\nabla w'=2\sum_kg'^kM^k,\qquad
a'_{ij}(x)=D_{h_j}A^i\bigl(g'(x)\bigr),
\]
each obtained from $g'$ by the same formula that produces $w$, $\nabla w$,
$a_{ij}$ from $g$.  At every point where $g=g'$ the two members of each pair
take the same value, so by $(1)$
\[
w=w',\qquad \nabla w=\nabla w',\qquad a_{ij}=a'_{ij}
\qquad\text{almost everywhere on }\Omega ,
\tag{1$'$}
\]
and likewise $w^\sigma=(w')^\sigma$ almost everywhere on $\Omega$ for every real
$\sigma$.
(The same predicate, unused in any other way here, is what Step 4 hands to
\noderef{DifferentiatedEquationWeak}.)  Next, both $w=\varepsilon+|g|^2$ and
$w'=\varepsilon+|g'|^2$ are $\ge\varepsilon>0$ everywhere, and $w'$ is
continuous on $\Omega$ by $(1)$; hence so is every real power $(w')^\sigma$,
because $\lambda\mapsto\lambda^\sigma$ is continuous on $(0,\infty)$.  Also
$|g|^2=w-\varepsilon\le w$, so by the Cauchy--Schwarz inequality (Preamble) in
the index $k$,
\[
|\nabla w|^2=4\Bigl|\sum_kg^kM^k\Bigr|^2
\le4\Bigl(\sum_k(g^k)^2\Bigr)\Bigl(\sum_k|M^k|^2\Bigr)
\le4\,w\sum_k|M^k|^2 ,
\qquad\text{i.e.}\qquad
\sum_k|M^k|^2\ \ge\ \frac{|\nabla w|^2}{4w} .
\tag{2}
\]

\emph{Step 3: the two chain-rule computations.}  Fix $k$ and define
$F_k,G\colon\R^n\to\R$ by
\[
F_k(y)=y_k\Bigl(\varepsilon+\sum_jy_j^2\Bigr)^{\frac t2},
\qquad
G(y)=\Bigl(\varepsilon+\sum_jy_j^2\Bigr)^{s} .
\]
Both are of class $C^1$ on $\R^n$: $y\mapsto\varepsilon+|y|^2$ is a polynomial,
hence smooth, and takes values in $[\varepsilon,\infty)\subseteq(0,\infty)$
because $\varepsilon>0$, while $\lambda\mapsto\lambda^{t/2}$ and
$\lambda\mapsto\lambda^{s}$ are smooth on $(0,\infty)$; a product of $C^1$
functions is $C^1$.  (This is the one place $\varepsilon>0$ is used for
smoothness: at $\varepsilon=0$ the map $y\mapsto|y|^{t}$ is not differentiable at
the origin for small $t>0$, which is why \S3 works at $\varepsilon>0$.)  Their
partial derivatives are
\[
D_jF_k(y)=\delta_{jk}\Bigl(\varepsilon+|y|^2\Bigr)^{\frac t2}
+t\,y_ky_j\Bigl(\varepsilon+|y|^2\Bigr)^{\frac t2-1},
\qquad
D_jG(y)=2s\,y_j\Bigl(\varepsilon+|y|^2\Bigr)^{s-1} .
\]
\emph{(3a) The primed family lies in $W^{1,2}_{\mathrm{loc}}(\Omega)$ with the
same gradients.}  The present node assumes
$g^j\in W^{1,2}_{\mathrm{loc}}(\Omega)$ with weak gradient $M^j$
(\noderef{MemW1pLocOn}); we record that $g'^j$ does too, with the \emph{same}
$M^j$.  By \noderef{MemW1pLocOn} this is two assertions, and each is an integral
one.  First the identity of \noderef{HasWeakGradientOn}: for every $\psi$ in
that node's test class and every $i$ we have
$\int_\Omega g^j\,D_i\psi\,dx=-\int_\Omega(M^j)_i\,\psi\,dx$; by $(1)$,
$g^j=g'^j$ almost everywhere on $\Omega$, so the two left-hand integrands agree
almost everywhere and the integrals are equal, giving
$\int_\Omega g'^j\,D_i\psi\,dx=-\int_\Omega(M^j)_i\,\psi\,dx$.  Second the
$L^2$ clauses: for compact $K\subseteq\Omega$ we have $g'^j=g^j$ almost
everywhere on $K$ and $M^j$ is not changed at all, so $g'^j\in L^2(K)$ and
$M^j\in L^2(K)$ follow from the corresponding statements for $g^j$.  Hence
\[
g'^j\in W^{1,2}_{\mathrm{loc}}(\Omega)\ \text{ with weak gradient }\ M^j
\qquad(j=1,\dots,n).
\tag{3a}
\]

\emph{(3b) The composition, performed on the representative.}  Apply
\noderef{WeakGradientVectorComp} with this $n$, with $m=n$, with $q=2$, with the
open set $\Omega$, with $G_j=g'^j$ and $DG_j=M^j$, and with $F=F_k$.  Its
hypotheses hold: $1\le n$ and $1\le2$; $\Omega$ is open; the
$W^{1,2}_{\mathrm{loc}}$ clause is $(3a)$; the family is locally uniformly
bounded, because for compact $K\subseteq\Omega$ the function $|g'|$ is
continuous on $K$ by $(1)$ and hence bounded there, say by $L$, and then
$|g'^j(x)|\le|g'(x)|\le L$ for \emph{every} $x\in K$ and every $j$; and $F_k$ is
$C^1$.  \emph{This hypothesis is the one place in the proof that forces the
representative on us: \noderef{WeakGradientVectorComp} asks for a bound at every
point of $K$, and the given $g$ need not satisfy one, since altering a weak
gradient on a null set leaves every hypothesis of the present node intact while
destroying any pointwise bound.}  Since $F_k(g'(x))=g'^k(x)\,w'(x)^{t/2}$, the
conclusion is
\[
\varphi'_k:=g'^k\,(w')^{\frac t2}\in W^{1,2}_{\mathrm{loc}}(\Omega)
\quad\text{with weak gradient}\quad
\Psi'_k=(w')^{\frac t2}M^k+\frac t2\,g'^k\,(w')^{\frac t2-1}\nabla w' ,
\]
where the displayed field is $\sum_jD_jF_k(g')\,M^j$ rewritten using
$\sum_jg'^jM^j=\tfrac12\nabla w'$.  The same node applied with $F=G$ in place of
$F_k$, under the same hypotheses, gives
\[
(w')^{s}\in W^{1,2}_{\mathrm{loc}}(\Omega)
\quad\text{with weak gradient}\quad
\sum_jD_jG(g')M^j=2s\,(w')^{s-1}\sum_jg'^jM^j=s\,(w')^{s-1}\nabla w' .
\]

\emph{(3c) Transport back to the unprimed quantities.}  By $(1')$,
$\varphi_k:=g^kw^{t/2}$ equals $\varphi'_k$ almost everywhere on $\Omega$, and
$\Psi_k:=w^{t/2}M^k+\tfrac t2g^kw^{t/2-1}\nabla w$ equals $\Psi'_k$ almost
everywhere on $\Omega$; likewise $w^s=(w')^s$ and
$s\,w^{s-1}\nabla w=s\,(w')^{s-1}\nabla w'$ almost everywhere on $\Omega$.  Now
repeat the unfolding of $(3a)$ verbatim with $(\varphi'_k,\Psi'_k)$ in place of
$(g^j,M^j)$: both sides of the \noderef{HasWeakGradientOn} identity and both
$L^2$ clauses of \noderef{MemW1pLocOn} are integrals, so each is unchanged when
its integrand is altered on a null set, and the membership passes to the
unprimed pair.  The same for $(w^s,s\,w^{s-1}\nabla w)$.  Therefore
\[
\varphi_k=g^k\,w^{\frac t2}\in W^{1,2}_{\mathrm{loc}}(\Omega)
\quad\text{with weak gradient}\quad
\Psi_k=w^{\frac t2}M^k+\frac t2\,g^k\,w^{\frac t2-1}\nabla w ,
\tag{3}
\]
and
\[
w^{s}\in W^{1,2}_{\mathrm{loc}}(\Omega)
\quad\text{with weak gradient}\quad
s\,w^{s-1}\nabla w ,
\tag{4}
\]
which is the identity $\nabla(w^s)=s\,w^{s-1}\nabla w$ the statement's second
display records, and which identifies the left-hand side of the assertion with
$s^2I$.

\emph{Step 4: $\varphi_k\zeta^2$ is an admissible test function for
(3.2) of the thesis.}  Put $K=\operatorname{supp}\zeta$, a compact subset of
$\Omega$, and
\[
\varphi=\varphi_k\,\zeta^2,
\qquad
g_\varphi=\zeta^2\Psi_k+2\zeta\,\varphi_k\,\nabla\zeta .
\]
We check that $\varphi\in W^{1,2}(\Omega)$ with weak gradient $g_\varphi$
(\noderef{MemW1pOn} at the exponent $2$) and that $\varphi$ has compact support
inside $\Omega$, which is exactly the test class of
\noderef{DifferentiatedEquationWeak}.

The support statement is immediate: $\zeta^2$ and
$\nabla(\zeta^2)=2\zeta\nabla\zeta$ vanish off $K$, so $\varphi$ and
$g_\varphi$ vanish off $K$, and $\operatorname{supp}\varphi\subseteq K\subseteq\Omega$
with $K$ compact.

The $W^{1,2}(\Omega)$ membership with weak gradient $g_\varphi$ is clause (2) of
\noderef{WeakGradientSmoothMul}, applied with the present dimension $n$, with
$q=2$, with the open set $\Omega$, with the pair $(\varphi_k,\Psi_k)$ in the
role of $(f,g)$, and with the smooth factor $\zeta^2$ in the role of $\psi$.
Its hypotheses hold: $1\le n$ and $1\le q$ because $2\le n$ and $q=2$;
$\Omega$ is open; $\varphi_k\in W^{1,2}_{\mathrm{loc}}(\Omega)$ with weak
gradient $\Psi_k$ is exactly $(3)$; and $\zeta^2$ is $C^\infty$ because $\zeta$
is, has compact support with
$\operatorname{supp}(\zeta^2)\subseteq\operatorname{supp}\zeta=K\subseteq\Omega$.
It is clause (2) and not clause (1) that applies, because the factor $\varphi_k$
is only locally in $W^{1,2}$ and because $\Omega$ is an arbitrary open set, of
possibly infinite volume; clause (2) assumes neither $\varphi_k\in
W^{1,2}(\Omega)$ nor $|\Omega|<\infty$.  The conclusion is that
$\varphi=\varphi_k\zeta^2\in W^{1,2}(\Omega)$ (\noderef{MemW1pOn}) with weak
gradient
\[
\zeta^2\Psi_k+\varphi_k\,\nabla(\zeta^2)
=\zeta^2\Psi_k+2\zeta\,\varphi_k\,\nabla\zeta
=g_\varphi ,
\]
the first equality being the classical product rule $D_i(\zeta^2)=2\zeta D_i\zeta$
for the smooth function $\zeta$.

\emph{Step 5: substitution and summation over $k$.}  Apply
\noderef{DifferentiatedEquationWeak} with the present
$n,p,\varepsilon,\gamma_0,\gamma_1,\Omega,A,u,g,Dg$, with the index $k$ and with
the test pair $(\varphi,g_\varphi)$ of Step 4.  Its hypotheses are $2\le n$,
$1<p$, $0<\varepsilon$, $0<\gamma_0$, $\gamma_0\le\gamma_1$ (Step 1), openness of
$\Omega$, the structure conditions (\noderef{EllipticStructureQ12}), the cited
regularity result (C1) (\noderef{CitedRegularityEpsPos}, a hypothesis of the
present node), the weak-solution property (\noderef{IsWeakSolutionOn}), the
membership of each $g^i$ in $W^{1,2}_{\mathrm{loc}}(\Omega)$ with weak gradient
$M^i$ (\noderef{MemW1pLocOn}), and the test-class conditions on
$(\varphi,g_\varphi)$, which are Step 4.  Its conclusion is
\[
\int_\Omega\sum_{i,j}a_{ij}\,(M^k)^j\,g_\varphi^i\,dx=0
\qquad(k=1,\dots,n).
\]
Insert $g_\varphi^i=\zeta^2\bigl(w^{t/2}(M^k)^i+\tfrac t2g^kw^{t/2-1}(\nabla w)_i\bigr)
+2\zeta g^kw^{t/2}(D_i\zeta)$, which is Step 4's field written out with $(3)$ and
$\varphi_k=g^kw^{t/2}$, and sum the $n$ identities.  Using
$\sum_kg^k(M^k)^j=\tfrac12(\nabla w)_j$ in the second and third groups of terms,
\[
\underbrace{\int_\Omega\zeta^2w^{\frac t2}\sum_k\sum_{i,j}a_{ij}(M^k)^j(M^k)^i}_{T_1}
+\underbrace{\frac t4\int_\Omega\zeta^2w^{\frac t2-1}\sum_{i,j}a_{ij}(\nabla w)_j(\nabla w)_i}_{T_2}
=-\underbrace{\int_\Omega\zeta\,w^{\frac t2}\sum_{i,j}a_{ij}(\nabla w)_j(D_i\zeta)}_{T_3}.
\tag{5}
\]
All three integrands vanish off $K$ and are integrable there: by (Q2) of the thesis,
$|a_{ij}|\le\gamma_1w^{p/2-1}$ almost everywhere on $\Omega$, and $w^{p/2-1}$,
$w^{t/2}$, $w^{t/2-1}$, $\zeta$ and $\nabla\zeta$ are bounded almost everywhere
on $K$ --- for the three powers of $w$, because $w'$ is continuous and
$\ge\varepsilon>0$ on the compact $K$, hence each $(w')^\sigma$ is bounded there,
and $w^\sigma=(w')^\sigma$ almost everywhere by $(1')$; a bound holding almost
everywhere dominates the same integral --- while $|M^k|\in L^2(K)$ for each
$k$ and $|\nabla w|\le2\sqrt w\,(\sum_k|M^k|^2)^{1/2}\in L^2(K)$ by $(2)$; so
each integrand is dominated by a constant multiple of a product of two
$L^2(K)$ functions, hence is in $L^1(K)$ by Cauchy--Schwarz (Preamble).  The
rearrangement of the finite sums and the splitting of the integrals is therefore
legitimate.

\emph{Step 6: the lower bound from (Q1) of the thesis.}  Condition (Q1) of the thesis at a
point $x$ where it holds, with $z=0$, $h=g(x)$ --- whose guard
$\varepsilon+|h|^2=w(x)>0$ is satisfied --- and an arbitrary $\xi$, reads
$\gamma_0w^{p/2-1}|\xi|^2\le\sum_{i,j}a_{ij}\xi^j\xi^i$.  Taking $\xi=M^k$ and
summing over $k$, and then taking $\xi=\nabla w$,
\[
\sum_k\sum_{i,j}a_{ij}(M^k)^j(M^k)^i\ \ge\ \gamma_0w^{\frac p2-1}\sum_k|M^k|^2
\ \ge\ \frac{\gamma_0}4\,w^{\frac p2-2}|\nabla w|^2 ,
\qquad
\sum_{i,j}a_{ij}(\nabla w)_j(\nabla w)_i\ \ge\ \gamma_0w^{\frac p2-1}|\nabla w|^2
\]
almost everywhere on $\Omega$, the middle step by $(2)$.  Multiplying by the
nonnegative weights $\zeta^2w^{t/2}$ and $\tfrac t4\zeta^2w^{t/2-1}$,
integrating (monotonicity of the integral, Preamble) and adding the exponents
--- $\tfrac t2+\tfrac p2-2=\tfrac{p+t}2-2$ and
$\tfrac t2-1+\tfrac p2-1=\tfrac{p+t}2-2$ ---
\[
T_1\ \ge\ \frac{\gamma_0}4\,I,
\qquad
T_2\ \ge\ \frac{\gamma_0t}4\,I,
\qquad\text{so}\qquad
\gamma_0\,\frac{1+t}4\,I\ \le\ T_1+T_2\ =\ -T_3\ \le\ |T_3| .
\tag{6}
\]
Here $I<\infty$: its integrand vanishes off $K$ and, by the bounds of Step 5, is
dominated on $K$ by a constant multiple of $|\nabla w|^2\in L^1(K)$.  Likewise
$J<\infty$.

\emph{Step 7: the upper bound from (Q2) of the thesis and Cauchy--Schwarz.}  By
(Q2) of the thesis, almost everywhere on $\Omega$,
\[
\Bigl|\sum_{i,j}a_{ij}(\nabla w)_j(D_i\zeta)\Bigr|
\le\gamma_1w^{\frac p2-1}\Bigl(\sum_j|(\nabla w)_j|\Bigr)\Bigl(\sum_i|D_i\zeta|\Bigr)
\le\gamma_1n\,w^{\frac p2-1}|\nabla w|\,|\nabla\zeta| ,
\]
the last step by Cauchy--Schwarz (Preamble) applied twice, to
$(|(\nabla w)_j|)_j$ against $(1,\dots,1)$ and to $(|D_i\zeta|)_i$ against
$(1,\dots,1)$, each giving a factor $\sqrt n$.  Hence, by the norm inequality
for the integral and monotonicity (Preamble),
\[
|T_3|\ \le\ \gamma_1n\int_\Omega w^{\frac{p+t}2-1}|\nabla w|\,|\nabla\zeta|\,|\zeta|\,dx
\ \le\ \gamma_1n\,I^{\frac12}J^{\frac12},
\]
where the second inequality is Cauchy--Schwarz (Preamble) for the two factors
\[
w^{\frac{p+t}4-1}|\nabla w|\,|\zeta|
\quad\text{and}\quad
w^{\frac{p+t}4}|\nabla\zeta| ,
\]
whose product is the integrand --- the exponents adding as
$\bigl(\tfrac{p+t}4-1\bigr)+\tfrac{p+t}4=\tfrac{p+t}2-1$ --- and whose squares
integrate to $I$ and $J$.  Combining with $(6)$,
\[
\gamma_0\,\frac{1+t}4\,I\ \le\ \gamma_1n\,I^{\frac12}J^{\frac12}.
\tag{7}
\]

\emph{Step 8: conclusion.}  If $I=0$ then the left-hand side of the assertion is
$s^2I=0$, while the right-hand side $cJ$ is nonnegative because $c>0$ and the
integrand of $J$ is nonnegative; so the assertion holds.  Assume $I>0$.  Then
$I<\infty$ by Step 6, so dividing $(7)$ by $I^{1/2}>0$ and squaring,
\[
I\ \le\ \Bigl[\frac{4\gamma_1n}{\gamma_0(1+t)}\Bigr]^2J .
\]
Multiplying by $s^2=\bigl(\tfrac{p+t}4\bigr)^2$,
\[
s^2I\ \le\ \Bigl(\frac{p+t}4\Bigr)^2\frac{16\,\gamma_1^2n^2}{\gamma_0^2(1+t)^2}\,J
= n^2\Bigl(\frac{\gamma_1}{\gamma_0}\Bigr)^2\Bigl(\frac{p+t}{1+t}\Bigr)^2J .
\]
Now $\gamma_1\le r\gamma_0$ with $\gamma_0>0$ gives $\gamma_1/\gamma_0\le r$, and
$\gamma_1/\gamma_0>0$ by Step 1, so $(\gamma_1/\gamma_0)^2\le r^2$; and for every
$t>0$
\[
\frac{p+t}{1+t}\le p ,
\]
because $p>1$ gives $p+t\le p+pt=p(1+t)$ and $1+t>0$.  Hence
\[
s^2I\ \le\ n^2r^2p^2\,J\ =\ c\,J ,
\]
which is the assertion.  This is the thesis' reduction ``since
$\frac{p+t}{1+4t}\le p$ for all $t>0$'' in the form the computation above
produces; the thesis' intermediate display carries the coefficient
$\gamma_0(\frac14+t)$ where the summation of Step 5 gives
$\gamma_0\frac{1+t}4$, and $\gamma_1n^2$ where (Q2) of the thesis gives $\gamma_1n$.
Neither difference affects (3.3) of the thesis: the constant is existential and both
numerical factors depend only on $n$, $p$ and the ratio $\gamma_1/\gamma_0$, so
the same final bound is reached with $c=n^2p^2r^2$ in place of the thesis'
unevaluated $c(n,p,\gamma_1/\gamma_0)$.
\end{proof}
